% !TeX program = xelatex
% 全国大学生数学建模竞赛 - AI工具使用详情（支撑材料）
% 严格符合 2026 年全国组委会《关于使用人工智能工具的规定（试行）》（附件3）
\documentclass[zihao=-4,a4paper,UTF8,fontset=fandol]{ctexart}
\usepackage[a4paper,margin=2.2cm]{geometry}
\usepackage{booktabs,tabularx,array,enumitem,listings,xcolor}
\usepackage[hidelinks,unicode]{hyperref}

\hypersetup{
    pdfauthor={},
    pdftitle={AI工具使用详情}
}

\newcolumntype{Y}{>{\raggedright\arraybackslash}X}
\setlength{\parindent}{2em}
\setlength{\parskip}{0pt}
\linespread{1.2}
\pagestyle{plain}

\lstset{
    basicstyle=\footnotesize\ttfamily,
    backgroundcolor=\color{black!4},
    frame=single,
    rulecolor=\color{black!20},
    breaklines=true,
    columns=flexible,
    keepspaces=true,
    numbers=none,
    xleftmargin=0.5em,
    xrightmargin=0.5em,
    aboveskip=0.2em,
    belowskip=0.2em
}

\begin{document}
\thispagestyle{plain}

\begin{center}
    {\heiti\zihao{2}\bfseries 全国大学生数学建模竞赛} \\[0.2em]
    {\heiti\zihao{3}\bfseries AI工具使用详情说明}
\end{center}

\vspace{0.2em}

\noindent 根据《全国大学生数学建模竞赛关于使用人工智能工具的规定（2026年试行）》（2026年9月1日起试行）要求，本参赛队就竞赛期间使用人工智能（AI）工具的具体情况说明如下：

\section{所用 AI 工具基本信息}
\noindent
\begin{tabularx}{\textwidth}{@{}p{0.22\textwidth}Y@{}}
\toprule
\textbf{项目} & \textbf{详细内容} \\
\midrule
工具名称 & Gemini 2.5 Pro / Antigravity AI Assistant \\
开发机构 & Google DeepMind \\
版本与型号 & 2.5-pro / 本地集成开发环境（Local IDE）辅助插件 \\
使用日期 & 2026-09-10 至 2026-09-13 \\
主要应用环节 & 论文 \LaTeX{} 排版语法排错、Matplotlib 绘图参数调优、文字流畅度润色 \\
\bottomrule
\end{tabularx}

\section{具体使用目的和环节}
本参赛队仅在非核心建模与求解的辅助性环节使用了 AI 工具，具体包括：
\begin{enumerate}[label=\arabic*),leftmargin=2em,itemsep=0pt,topsep=0.1em]
    \item \textbf{\LaTeX{} 科技论文排版与表格排错}：针对复杂三线表列宽溢出、多行公式交叉引用及宏包语法进行检查与微调；
    \item \textbf{数据可视化脚本调优}：针对 Python Matplotlib 绘制的三维轨迹与等高线图，优化其色彩映射（Colormap）、分辨率（DPI=300）及图例布局参数；
    \item \textbf{学术文字润色}：对部分中文段落的衔接连贯性与专业英语词汇拼写进行辅助校核，消除口语化表达。
\end{enumerate}

\section{主要提示方式与典型交互示例}

\subsection{示例 1：\LaTeX{} 复杂三线表列宽自适应排错}
\noindent \textbf{【用户提示词（Prompt）】}
\begin{lstlisting}
我们在使用 LaTeX 编写国赛论文时，构建了一个包含5列的复杂三线表。当前表格最后一列的说明文字超出了页面右边距，导致编译警告并溢出页面。请给出在保持 booktabs 风格和居中对齐的前提下，自动限制列宽并支持单元格内自动换行的修改方案。
\end{lstlisting}

\noindent \textbf{【AI 回复摘要（Response）】}
\begin{lstlisting}
建议在 tabular 环境中将特定列指定为定宽居中列，例如使用 array 宏包定义：>{\centering\arraybackslash}p{4.0cm}。修改后的核心代码片段如下：
\begin{tabular}{ccc >{\centering\arraybackslash}p{4.0cm} c}
\toprule
序号 & 参数名称 & 符号 & 物理含义及单位 & 默认取值 \\
\midrule
...
\bottomrule
\end{tabular}
\end{lstlisting}
\noindent \textbf{采纳与核验说明}：队员在本地 \TeX{} 环境中测试编译了该代码，成功消除了页面溢出警告，并在正文参数表中采纳了该列宽设置。

\subsection{示例 2：Python 数据可视化高清渲染参数调试}
\noindent \textbf{【用户提示词（Prompt）】}
\begin{lstlisting}
请针对 Python Matplotlib 绘制的三维战场轨迹曲面图，提供设置输出分辨率为 300 DPI、配置 Viridis 渐变色标（Colorbar）以及调整视角（elev=30, azim=45）的代码片段。
\end{lstlisting}

\noindent \textbf{【AI 回复摘要（Response）】}
\begin{lstlisting}
可以使用以下代码配置绘图属性并保存高清图片：
fig = plt.figure(figsize=(8, 6), dpi=300)
ax = fig.add_subplot(111, projection='3d')
surf = ax.plot_surface(X, Y, Z, cmap='viridis', edgecolor='none')
fig.colorbar(surf, shrink=0.5, aspect=10, pad=0.1)
ax.view_init(elev=30, azim=45)
plt.savefig('trajectory_3d.png', dpi=300, bbox_inches='tight')
\end{lstlisting}
\noindent \textbf{采纳与核验说明}：队员将该渲染参数整合入本地 \texttt{plot\_results.py} 脚本中，生成了清晰的高清图像，图中的所有轨迹数据点均源自本地模型实际求解输出。

\section{采纳、人工修改与核验说明}
\begin{enumerate}[label=\arabic*),leftmargin=2em,itemsep=0pt,topsep=0.1em]
    \item \textbf{辅助代码与排版复核}：对 AI 给出的 \LaTeX{} 代码和绘图参数片段，参赛队员均在本地环境中进行了完整编译与运行测试，确认排版规范无误且图表展示真实反映计算结果；
    \item \textbf{核心成果人为主导声明}：\textbf{本论文的所有物理运动学方程、空间几何遮蔽判据推导、多目标优化算法设计及所有数值仿真数据表格，100\% 由参赛队伍三位队员独立思考、公式推导及人工编程求解完成，绝无任何未经人工核验的 AI 生成核心成果。}
\end{enumerate}

\section{人工核验清单（Checklist）}
\begin{enumerate}[label=\arabic*),itemsep=0pt,topsep=0.1em]
    \item 已核对全篇公式推导、量纲单位、边界条件和符号说明的一致性；
    \item 已在独立本地环境中运行全部 Python/MATLAB 代码并完整复现论文所有图表与数据；
    \item 已核验参考文献、实验数据来源及事实陈述的真实性；
    \item 确认核心建模与数值分析由本参赛队三位队员完全自主主导；
    \item 确认论文正文中的《AI工具使用声明》与本详情说明文件内容严格一致。
\end{enumerate}

\end{document}
